\documentclass[10pt,a4paper]{article}

\usepackage[utf8]{inputenc}
\usepackage[T1]{fontenc}
\setlength{\textwidth}{6.5in}
\setlength{\oddsidemargin}{0in}
\setlength{\evensidemargin}{0in}
\setlength{\textheight}{9in}
\setlength{\topmargin}{-0.5in}
\usepackage{amsmath,amssymb}
\usepackage{array}
\usepackage{booktabs}
\usepackage{cite}
\usepackage{hyperref}
\usepackage{url}

\hypersetup{
  colorlinks=true,
  linkcolor=black,
  citecolor=black,
  urlcolor=blue,
  breaklinks=true
}

\title{Reproducible Hybrid Recommendation for Vietnamese Retail}
\author{}
\date{}

\begin{document}
\maketitle

\begin{abstract}
Offline recommender results are difficult to interpret when model identity is
reported without the data, split, candidate, masking, objective, and evaluator
conditions that produced the result. This paper presents a provenance-aware
protocol for assessing a decoupled Wide-and-Deep two-tower recommender in a
controlled Vietnamese retail setting. The protocol fixes a temporal
novel-purchase task, full-catalog ranking, seen-item masking, deterministic tie
handling, per-user metrics, and a hierarchical uncertainty plan. It separates
source-native reproduction, public protocol validation, controlled-benchmark
comparison, and external validation into distinct evidence namespaces. The
completed public lane contains one deterministic Pop run and three BPR runs on
GroupLens MovieLens 100K; these rows are reported as descriptive protocol
evidence only. The proposed Hybrid has not been admitted to the paper result
set, and the controlled benchmark comparison remains a follow-up experiment.
The contribution is therefore a reproducible comparison design with explicit
evidence boundaries, rather than a conclusion about a settled model ordering.
\end{abstract}

\noindent\textbf{Keywords:} recommender systems, reproducibility, full-catalog evaluation, two-tower models, association rules, cold-item recommendation

% Stage 2 manuscript fragment. Included by master_draft_stage2.tex.
% The adjacent ref/anchor comments are provenance markers for the ARS audit.

\section{Introduction}

Retail recommendation is often presented as a model-selection problem, but an
offline result is first a statement about a defined ranking task. A system must
decide which user history is available, which items are eligible, what counts as
a relevant future event, and how scores are converted into an ordered list. In
large-scale services, candidate generation and ranking are also separate system
roles rather than interchangeable names for one predictor
\cite{covington2016_youtube} %<!--ref:covington2016_youtube--><!--anchor:section:Sections 2 System Overview, 3 Candidate Generation, and 4 Ranking-->
\cite{yi2019_ndr} %<!--ref:yi2019_ndr--><!--anchor:section:Sections 2.3 Two-tower Models, 3 Modeling Framework, and 5 Neural Retrieval System-->.
This distinction matters in retail settings with sparse user--item observations,
changing purchase sequences, item metadata, and products with little or no
training history. The same model label can describe different estimands when
the candidate universe or target event changes.

The present study addresses that problem in a controlled Vietnamese retail
setting. The benchmark is intended to represent the data structures needed for
retail ranking, including user histories, product content, categories, price
signals, and purchase or basket events. Its behavior is controlled and
semi-synthetic rather than a direct sample of Vietnamese shoppers. That choice
allows the protocol to expose specific mechanisms, such as rule-aligned
co-purchase signals and item-side content, while limiting the claims that can be
made about natural behavior. The benchmark declaration specifies 5,000 users,
5,200 items, 823,371 interactions, and a designated item-cold subset; these are
manifest-level design facts and remain subject to the lineage receipt for the
follow-up experiment.

\subsection{Why comparability is a research problem}

Offline recommender comparisons are sensitive to choices that are sometimes
reported as implementation details. Temporal partitioning changes the
information available at prediction time, while target construction determines
which events enter the evaluation population. A recent study of splitting
strategies demonstrates that sequential-recommender outcomes can vary with the
split policy
\cite{gusak2025_time_split} %<!--ref:gusak2025_time_split--><!--anchor:section:Figure 1; Sections 1–3; global temporal split and results subsections-->
and a systematic replicability analysis of BERT4Rec shows that training and
implementation decisions are part of the reproducibility question
\cite{petrov2022_a_systematic_review_and_replicability_study_of_bert4rec_fo} %<!--ref:petrov2022_a_systematic_review_and_replicability_study_of_bert4rec_fo--><!--anchor:section:Implementation review and RQ2 training-time analysis, including Table 4-->.
Candidate construction, seen-item masking, negative sampling, cutoff, tie
handling, and aggregation consequently define the quantity being estimated.
Metrics with the same name are not automatically comparable when those inputs
differ.

This paper responds by treating the evaluation protocol as a first-class study
artifact. The protocol freezes an immutable temporal split, a complete item
catalog, eligibility rules, relevance semantics, masking, deterministic ranking,
and metric implementation. Each model exports scores to an independent shared
evaluator. The evaluator, rather than the model implementation, applies
masking, ranking, metric calculation, and per-user aggregation. This separation
does not make every method equivalent. It makes the remaining differences
inspectable and prevents a source-native number from silently becoming a result
under a different task contract.

\subsection{Distinct signals and the conditional Hybrid question}

The candidate design combines a deep two-tower scorer with a rule-derived Wide
component. The rationale is a complementarity hypothesis: a rule branch may
retain explicit co-purchase regularities, while a content-aware deep branch may
generalize beyond interactions that have already been observed. This hypothesis
is related to, but not identical with, the Wide \& Deep decomposition of linear
cross-product memorization and embedding-based generalization
\cite{cheng2016_wide_deep} %<!--ref:cheng2016_wide_deep--><!--anchor:section:Sections 2 Recommender System Overview and 3.1-3.3 Wide, Deep, and Joint Training-->.
It is also distinct from item-based collaborative filtering, which aggregates
neighbor evidence from interaction overlap
\cite{sarwar2001_itemcf} %<!--ref:sarwar2001_itemcf--><!--anchor:section:Sections 3.1 Item Similarity Computation and 3.2 Prediction Computation-->,
and from Bayesian Personalized Ranking, which specifies a pairwise objective
for implicit feedback
\cite{rendle2009_bpr} %<!--ref:rendle2009_bpr--><!--anchor:section:Sections 4.1 BPR Optimization Criterion, 4.2 BPR Learning Algorithm, and 4.3.1 Matrix Factorization-->.
The distinction is substantive: a change in encoder, objective, negative
distribution, or system stage can change a result independently of the other
components.

Association rules provide a mechanism vocabulary, not an automatic ranking
claim. Apriori defines support and confidence as thresholds for mining frequent
item relationships
\cite{agrawal1994_apriori} %<!--ref:agrawal1994_apriori--><!--anchor:section:Introduction, formal statement of the association-rule problem and minimum support/confidence thresholds-->,
and multi-item rule systems show how such relationships can be operationalized
for recommendation under a particular dataset and evaluation design
\cite{ghoshal2014_multi_item_rules} %<!--ref:ghoshal2014_multi_item_rules--><!--anchor:section:Abstract; Section 2 problem statement; Section 6 comparison with collaborative filtering and matrix factorization-->.
Prior hybrid and next-basket studies provide useful component-comparison
precedents, including combinations of preference, popularity, transition, and
association signals
\cite{liu2009_hybrid_seq_cf} %<!--ref:liu2009_hybrid_seq_cf--><!--anchor:section:Replacement work Sections 3 and 3.1; Figures 1–4; Conclusions-->
\cite{peng2022_ham} %<!--ref:peng2022_ham--><!--anchor:section:Modeling sections for long-term preference and association components; Section 6.6 ablation study and Table 13-->
\cite{peng2023_m2} %<!--ref:peng2023_m2--><!--anchor:section:Section 4.1 model components; Section 6.5 ablation study-->,
but those precedents do not determine the contribution of the proposed branch
under the locked v5 cohort.

\subsection{Research question and contributions}

The primary research question is: under a fixed temporal novel-purchase
full-catalog protocol on the controlled Vietnamese retail benchmark v5, does
the proposed decoupled Wide-and-Deep Two-Tower Hybrid improve ranking
effectiveness relative to the strongest faithfully reproduced baseline selected
using validation data? Three subsidiary questions concern item-side cold
performance, the incremental role of the rule branch, and replication on a
compatible public commerce dataset. A fourth concerns efficiency, but it is
opened only after the accuracy artifacts have passed their evidence gates.

The paper makes three contributions that are supported at this stage. First, it
specifies a provenance-aware protocol in which data, split, candidate catalog,
masking, evaluator, seeds, and model artifacts are explicit. Second, it separates
source-native reproduction, public protocol validation, controlled-v5 comparison,
and external validation into distinct evidence namespaces. Third, it reports an
audited public protocol validation lane and preserves its limitations rather
than using it to stand in for the unexecuted Hybrid comparison. The resulting
claim is therefore about reproducible evaluation design and bounded evidence,
not about a settled ordering of recommender architectures. The controlled
Hybrid experiment remains a separately audited follow-up task.

The remainder of the paper first synthesizes prior work by signal and evaluation
contract. It then defines the data and ranking protocol, describes the evidence
lanes, reports the permitted public validation rows, and discusses the limits of
the current evidence.

\section{Related Work}

\subsection{Evaluation and reproducibility}

Recommendation literature often compares model families through headline
metrics, yet the metric is only interpretable together with the prediction
event and candidate set. The split-sensitivity evidence in sequential
recommendation makes temporal construction a substantive design choice
\cite{gusak2025_time_split} %<!--ref:gusak2025_time_split--><!--anchor:section:Figure 1; Sections 1–3; global temporal split and results subsections-->,
while the BERT4Rec replicability study identifies implementation and training
choices that affect whether a reported result can be reproduced
\cite{petrov2022_a_systematic_review_and_replicability_study_of_bert4rec_fo} %<!--ref:petrov2022_a_systematic_review_and_replicability_study_of_bert4rec_fo--><!--anchor:section:Implementation review and RQ2 training-time analysis, including Table 4-->.
These observations support a narrow methodological conclusion: a comparison
should freeze the target event, temporal boundary, candidate universe, masking,
metric code, aggregation, and training budget before interpreting model
differences.

The present protocol consequently distinguishes source-native evidence from
harmonized evidence. A source-native result retains the source implementation,
data preparation, split, candidate policy, and evaluator. A harmonized result
uses a common data and evaluation contract only after the adapter's deviations
are declared. The distinction protects against a common category error in which
a number reproduced from a paper is treated as if it had been generated under a
new dataset or evaluator. It also makes a failed or incomparable reproduction a
visible outcome rather than a reason to substitute a convenient surrogate.

The positioning is thus deliberately limited. Prior work motivates protocol
control, but it does not supply a ranking for the v5 benchmark or authorize
changing its frozen split after observing outcomes.

\subsection{Collaborative, neural, and sequence recommendation}

Item-based collaborative filtering estimates item similarity from historical
co-interaction and aggregates neighboring items for a user
\cite{sarwar2001_itemcf} %<!--ref:sarwar2001_itemcf--><!--anchor:section:Sections 3.1 Item Similarity Computation and 3.2 Prediction Computation-->.
BPR instead defines a pairwise preference objective for implicit feedback and
can be instantiated with matrix factorization or another differentiable
scorer
\cite{rendle2009_bpr} %<!--ref:rendle2009_bpr--><!--anchor:section:Sections 4.1 BPR Optimization Criterion, 4.2 BPR Learning Algorithm, and 4.3.1 Matrix Factorization-->.
These methods may share user and item identifiers, but they encode different
assumptions about similarity, optimization, and unobserved interactions.

Neural collaborative filtering learns nonlinear interaction functions from user
and item identifiers
\cite{he2017_ncf} %<!--ref:he2017_ncf--><!--anchor:section:Sections 3.1 General Framework, 3.2 GMF, 3.3 MLP, and 3.4 Fusion of GMF and MLP-->.
DeepFM, by contrast, combines low- and high-order interactions over sparse
feature fields for a click-through-rate task
\cite{guo2017_deepfm} %<!--ref:guo2017_deepfm--><!--anchor:section:Section 2 Our Approach, especially 2.1 DeepFM-->.
The shared presence of embeddings does not make their input or target contracts
interchangeable. This distinction informs the adapter registry: a model keeps
its native objective and feature assumptions, while its score output is mapped
to the common evaluator.

Sequential models add another source of non-equivalence. SASRec uses causal
self-attention for next-item prediction
\cite{kang2018_sasrec} %<!--ref:kang2018_sasrec--><!--anchor:section:Section III-C causality; Sections III-E and III-F prediction/objective; Section IV-D evaluation metrics-->,
whereas BERT4Rec uses bidirectional sequence encoding and a masked-item Cloze
objective
\cite{sun2019_bert4rec} %<!--ref:sun2019_bert4rec--><!--anchor:section:Section 3.3 Cloze task; Section 4.2 task settings and evaluation metrics-->.
Their source evaluations use paper-native task and candidate designs. A next-item
objective is therefore not identical to a next-basket set objective, and a
source metric cannot be copied into the v5 table without a protocol bridge.

These families supply comparator and objective precedents, not directly
comparable v5 results.

\subsection{Wide and Deep, Two-Tower, and semantic representations}

Wide \& Deep jointly trains an explicit linear cross-product component for
memorization with an embedding and multilayer component for generalization
\cite{cheng2016_wide_deep} %<!--ref:cheng2016_wide_deep--><!--anchor:section:Sections 2 Recommender System Overview and 3.1-3.3 Wide, Deep, and Joint Training-->.
That decomposition is useful for stating a mechanism hypothesis, but it does not
identify the Wide component with association-rule mining. In the present design,
the rule features are fit from the training transactions and exported as a
separate score. The no-Wide condition is required to assess whether the branch
contributes under the same evaluator.

Two-tower systems are frequently used in large-corpus retrieval. The YouTube
system separates candidate generation and ranking
\cite{covington2016_youtube} %<!--ref:covington2016_youtube--><!--anchor:section:Sections 2 System Overview, 3 Candidate Generation, and 4 Ranking-->,
and neural retrieval work shows that negative-distribution and sampling-bias
choices can be part of the model contract
\cite{yi2019_ndr} %<!--ref:yi2019_ndr--><!--anchor:section:Sections 2.3 Two-tower Models, 3 Modeling Framework, and 5 Neural Retrieval System-->.
DirectAU further illustrates that an objective can target alignment and
uniformity rather than the pairwise objective used by BPR
\cite{wang2022_directau} %<!--ref:wang2022_directau--><!--anchor:section:Sections 2.2 Alignment and Uniformity, 3 Alignment and Uniformity in CF, and 4 Direct Optimization-->.
The protocol must consequently identify whether a tower is a retriever,
pre-ranker, or full-catalog ranker and must record its negative sampler.

Recent tower designs make the same boundary visible from another direction.
ContextGNN introduces pair-aware local graph context while retaining an
exploratory two-tower fallback
\cite{yuan2025_contextgnn} %<!--ref:yuan2025_contextgnn--><!--anchor:section:Abstract and Method sections describing pair-wise familiar-item representations and the two-tower exploratory fallback-->,
and T2Diff adds cross-interaction during training while targeting efficient
matching
\cite{wang2025_t2diff} %<!--ref:wang2025_t2diff--><!--anchor:section:Introduction and Sections 4.2 Diffusion-based Cross Interaction, 4.3 Mixed-Attention Two-Tower, and 4.4 Optimization-->.
These works are useful controls for naming the pair-agnostic boundary. They do
not imply an architecture ordering under the v5 full-catalog contract.

The positioning is therefore specific: this paper evaluates a decoupled
architecture under one locked ranking contract and keeps encoder complexity,
objective, sampling, and system stage as separate explanatory variables.

\subsection{Apriori, association rules, and next-basket recommendation}

Apriori defines a transaction-frequency support threshold and a conditional
confidence threshold for mining association rules
\cite{agrawal1994_apriori} %<!--ref:agrawal1994_apriori--><!--anchor:section:Introduction, formal statement of the association-rule problem and minimum support/confidence thresholds-->.
Multi-item recommendation work demonstrates that such rules can be used to
construct a recommendation list, while its comparisons remain conditional on
the source data and protocol
\cite{ghoshal2014_multi_item_rules} %<!--ref:ghoshal2014_multi_item_rules--><!--anchor:section:Abstract; Section 2 problem statement; Section 6 comparison with collaborative filtering and matrix factorization-->.
The distinction is important for the proposed Wide branch: a mined rule is a
feature or score source, not evidence that the full ranking task has been
solved.

Next-basket research also warns against treating all basket events as one
behavioral regime. A reality-check analysis uses repeat ratios and baseline
behavior to expose how aggregate accuracy can be shaped by repetition
\cite{li2023_nbr_reality} %<!--ref:li2023_nbr_reality--><!--anchor:section:Figure 1 and Introduction contribution bullets; repeat-ratio and baseline analyses-->,
and exploration-focused work isolates the harder question of recommending novel
items
\cite{li2023_repetition_exploration} %<!--ref:li2023_repetition_exploration--><!--anchor:section:Introduction definitions and research questions; exploration-only analyses-->.
Mask-Swap and its BTBR baseline offer a direct next-novel-basket formulation,
but their reported conditions do not warrant a universal conclusion across
datasets or model families
\cite{li2023_mask_swap} %<!--ref:li2023_mask_swap--><!--anchor:section:Sections 4.1 and 4.2 mask/swap strategy; Section 5 BTBR model; Section 6 experiments-->.
Repeat/explore-aware graph work supplies further component and ablation
precedent under its own data and protocol
\cite{mansouri2026_repeat_explore_lightgcn} %<!--ref:mansouri2026_repeat_explore_lightgcn--><!--anchor:section:Abstract; Table 5 repeat/explore results; Section 6 ablation study; limitations-->.

These studies motivate a train-defined rule-aligned cohort and a no-Wide
ablation. They do not estimate the present mechanism question, and aggregate
basket accuracy alone does not establish a novel-item gain.

\subsection{Graph and contrastive recommendation}

LightGCN simplifies graph convolution by propagating trainable user and item
embeddings over observed interaction edges
\cite{he2020_lightgcn} %<!--ref:he2020_lightgcn--><!--anchor:section:Sections 3.1.1 Light Graph Convolution, 3.1.2 Layer Combination, and 3.3 Model Training-->.
Graph-contrastive controls such as SimGCL and LightGCL add contrastive or
global-relation objectives within their own observed-graph settings
\cite{yu2022_simgcl} %<!--ref:yu2022_simgcl--><!--anchor:section:Sections 2 Empirical Investigation and 3 SimGCL-->
\cite{cai2023_lightgcl} %<!--ref:cai2023_lightgcl--><!--anchor:section:Method section on SVD-guided global collaborative relation modeling and contrastive learning; Experiments-->.
An improvement on a sparse observed graph does not establish a representation
for an item with zero training edges. The v5 protocol therefore reports sparse
and strict zero-edge cohorts separately and does not infer cold-item capability
from a graph score alone.

Graph results motivate comparator coverage and precise cold-item definitions,
but they do not establish v5 performance or zero-edge performance.

\subsection{Content, cold-item, and transfer learning}

Cold-item analysis begins with the training graph. A strict zero-edge item has no
training interaction; an item with one or more training edges may be rare, but it
belongs to a different sparse cohort. DropoutNet formulates cold-start training
through an item-side framework
\cite{volkovs2017_dropoutnet} %<!--ref:volkovs2017_dropoutnet--><!--anchor:section:Section 2, Framework; Section 4.2, Training for Cold Start; Section 5.1, CiteULike-->,
and ALDI defines its own cold-item distillation setting
\cite{huang2023_aldi} %<!--ref:huang2023_aldi--><!--anchor:section:Section 2.1, Problem Formulation; Section 2.2, Existing Solutions-->.
Their thresholds and denominators cannot simply be transferred to v5. A
reproducible cold-item analysis must construct cohorts before scoring and report
warm, sparse, and zero-edge denominators separately.

Text encoders provide a useful item-side representation, but representation
quality is not a ranking outcome. Sentence-BERT supplies a rationale for
semantic sentence embeddings
\cite{reimers2019_sbert} %<!--ref:reimers2019_sbert--><!--anchor:section:Abstract; Section 1, Introduction; Section 3, Training Details-->,
whereas AlphaRec evaluates language representations inside a particular
recommendation and zero-shot protocol
\cite{sheng2025_alpharec} %<!--ref:sheng2025_alpharec--><!--anchor:section:Section 4.1, AlphaRec; Section 5.2, Zero-Shot Recommendation; Section 6, Limitations-->.
Sparse content-based cold-item work likewise defines its own problem and
evaluation setting
\cite{meehan2026_semco} %<!--ref:meehan2026_semco--><!--anchor:section:Section 1, Introduction; Section 3.1, Problem Definition; Sections 3.2-3.3-->.
The proposed pipeline therefore requires a recommender objective and a shared
task-specific evaluator in addition to an embedding recipe.

Transferable item representations retain target-domain assumptions. UniSRec
includes downstream adaptation
\cite{hou2022_unisrec} %<!--ref:hou2022_unisrec--><!--anchor:section:Figure 1; Sections 2.2-2.4, especially 2.4 Adaptation to Downstream Recommendation Tasks-->,
VQ-Rec evaluates vector-quantized transfer under its experimental setup
\cite{hou2023_vqrec} %<!--ref:hou2023_vqrec--><!--anchor:section:Sections 2.1-2.3; Section 3.1 Experimental Setup-->,
and UTGRec includes downstream fine-tuning
\cite{zheng2026_utgrec} %<!--ref:zheng2026_utgrec--><!--anchor:section:Section 2.4.2, Downstream Fine-tuning; Sections 3.1.1 and 3.1.3-3.1.4; Section 3.3, Ablation Study-->.
An architecture inspired by these works is not a replication unless the encoder,
data lineage, adaptation procedure, split, candidates, metrics, and statistics
are matched or every deviation is declared. Shared platform or benchmark
lineage also means that pretraining and downstream results cannot be counted as
independent replications merely because their domain labels differ
\cite{meehan2025_cold_popbias} %<!--ref:meehan2025_cold_popbias--><!--anchor:section:Section 1, Introduction; Sections 3.1-3.3-->
\cite{meehan2026_semco} %<!--ref:meehan2026_semco--><!--anchor:section:Section 1, Introduction; Section 3.1, Problem Definition; Sections 3.2-3.3-->.
Content aligned to collaborative teachers can also inherit popularity structure,
so content availability alone does not guarantee an unbiased or effective
cold-item ranking
\cite{meehan2025_cold_popbias} %<!--ref:meehan2025_cold_popbias--><!--anchor:section:Section 1, Introduction; Sections 3.1-3.3-->.

The positioning is conditional on full contract alignment. External and
cold-item conclusions require matched encoder, data, adaptation, split,
candidate set, metric, and statistical procedure; a source architecture name is
not sufficient evidence.

\section{Methodology and Protocol}

\subsection{Study design and estimand}

The study is designed as a quantitative comparative offline benchmark. Its unit
of analysis is a per-user metric vector nested within a training seed. This
choice preserves the dependence created when two models are evaluated for the
same users and permits uncertainty to be reported at both the user and seed
levels. The estimand is comparative algorithmic performance under a declared
dataset and evaluation protocol. It is not an estimate of changes in user
behavior, revenue, satisfaction, or long-term production outcomes.

The central comparison is defined before the follow-up test is opened. The
proposed model produces a score for every candidate item, and a locked baseline
registry supplies the comparison score. The strongest baseline is selected from
validation means using a fixed rule and is then frozen for the test evaluation.
This sequence prevents test outcomes from determining the comparator, features,
or stopping rule. A null or negative contrast remains a valid answer to the
research question.

The protocol separates three operations that are often conflated. Training
defines model parameters using only the training portion. Score export produces
an ordered-catalog-compatible score artifact without calculating final metrics.
Evaluation applies masking, ranking, relevance, and aggregation independently.
This separation makes the score-producing implementation auditable and keeps
metric semantics identical across model families.

\subsection{Dataset and temporal protocol}

The controlled v5 benchmark is a Vietnamese retail data construction with
semi-synthetic behavior. Its manifest declares user histories, item content,
category and price signals, and purchase or basket events. It also declares
5,000 users, 5,200 items, 823,371 interactions, and a designated item-cold
subset. These values are treated as dataset-manifest facts until the fresh
follow-up snapshot passes raw-to-internal reconciliation; they are not model
outcomes.

The declared temporal boundaries are train from 2026-01-01 through 2026-06-19,
validation from 2026-06-20 through 2026-07-10, and test from 2026-07-11
through 2026-08-01, all in UTC. User and item identifiers are mapped through an
immutable raw-to-internal table. The adapter checks schema, duplicate rows,
malformed records, timestamp ordering, user and item coverage, item-text
provenance, and the availability of basket or co-purchase structure. An
exclusion is recorded with its count and reason rather than silently removed.

The snapshot is admitted as a comparison input only when these checks form a
continuous chain. The raw provider bytes, the normalized tables, the
raw-to-internal identifier map, and the split manifest must each have a
content hash. The split hash must resolve to the protocol version used by the
runner, and the protocol version must resolve to the evaluator version used for
the reported scores. This conjunctive rule matters because a correct model
can still produce an uninterpretable result when, for example, an identifier
map is regenerated after the split or item text is read from a later catalog
revision. A failed link is recorded as an incomplete or diagnostic input; it
is not repaired by silently substituting a nearby snapshot. Dataset
suitability and model effectiveness are therefore assessed as separate
questions. The former concerns whether the data can support the declared
estimand, whereas the latter concerns the behavior of a model after the
estimand has been fixed.

The public protocol validation uses GroupLens MovieLens 100K as a separate
namespace and is not a substitute for the controlled v5 snapshot. It uses a
locked RecBole 1.2.1 environment and a full-sort evaluation path. The public
lane is useful for checking that a source-bound pipeline can materialize data,
train a reference method, export metrics, and preserve configuration and
checkpoint bindings. It does not establish that the same pipeline answers the
Vietnamese retail research question.

\subsection{Candidates, relevance, masking, and cohorts}

For each eligible user $u$, the candidate set is the complete item catalog after
removing items already observed in the allowed history:
\begin{equation}
  C_u = I \setminus H_u^{\mathrm{seen}}.
\end{equation}
The catalog order is frozen by raw product identifier. Scores are ranked in
descending order, with raw identifier ascending as the deterministic tie break.
No sampled candidate set is used when full-catalog scoring is available. This
policy makes the candidate universe explicit and avoids interpreting a sampled
metric as an exact-catalog quantity.

The primary target is a novel organic purchase in the target split. A user is
eligible only when at least one such positive event remains after the protocol
filters. Binary relevance is used for the top-$K$ metrics, with $K=10$. For a
user with relevant set $T_u$, the normalized discounted cumulative gain is
\begin{equation}
  \mathrm{NDCG@10}_u = \frac{\mathrm{DCG@10}_u}
  {\mathrm{IDCG@10}_u},
\end{equation}
where the ideal denominator uses $\min(|T_u|,10)$. Hit rate is the indicator
that at least one relevant item appears in the first ten positions. Recall is
reported separately because its denominator is the number of relevant target
items, not a hit indicator. Macro per-user GAUC is computed as the exact AUC
between each user's relevant positives and unseen non-positives, with average
rank for ties, and then averaged across eligible users.

The item-cold analysis is defined from the training graph. A strict zero-edge
item has no training interaction; an item with one or more training edges is
reported, if needed, in a sparse cohort rather than the zero-edge cohort. Warm,
sparse, and zero-edge denominators are constructed before scoring. The rule-
aligned cohort is likewise constructed from training transactions only. These
cohort definitions prevent a blended denominator from being presented as a
strict cold-item result.

\subsection{Models and adapter provenance}

The registry begins with Random and MostPop sanity controls, followed by
faithfully adapted reference methods where source and environment bindings are
available. ItemCF and BPR represent different collaborative contracts: the
former aggregates item-neighbor evidence, while the latter optimizes pairwise
implicit preferences \cite{sarwar2001_itemcf} %<!--ref:sarwar2001_itemcf--><!--anchor:section:Sections 3.1 Item Similarity Computation and 3.2 Prediction Computation-->
\cite{rendle2009_bpr} %<!--ref:rendle2009_bpr--><!--anchor:section:Sections 4.1 BPR Optimization Criterion, 4.2 BPR Learning Algorithm, and 4.3.1 Matrix Factorization-->.
The reference adapter retains the native objective and records repository
revision, dataset mapping, split, sampler, configuration, checkpoint rule, and
score-export revision. A source-native result is not relabeled as a harmonized
result merely because its scores can be read by the shared evaluator.

The proposed scope contains three related conditions. \texttt{IndependentDeepTwoTower}
uses user-history and item-content representations without the rule branch.
\texttt{WideRuleOnly} uses association features fit from training transactions.
\texttt{ProposedHybrid} combines the deep score and the rule score after the
normalization and additive fusion specified in a frozen configuration. All three
conditions emit component or total scores through the same \texttt{ScoreArtifact}
contract. The association rules are fit before validation and test data are
accessible, and the negative sampler draws only from items absent from the
user's training history.

The model descriptor also records the feature set, native objective, negative
sampler, implementation provenance, configuration hash, adapter revision, and
checkpoint hash. A text encoder is an input representation, not an outcome;
recommendation quality is assessed only after the recommender objective and
ranking evaluator have been applied. This constraint follows the distinction
between semantic embedding quality and task-specific recommendation evaluation
in prior work \cite{reimers2019_sbert} %<!--ref:reimers2019_sbert--><!--anchor:section:Abstract; Section 1, Introduction; Section 3, Training Details-->.

\subsection{Metrics and statistical plan}

NDCG@10 is the primary outcome. HR@10 and macro per-user GAUC are supporting
outcomes, and Recall@10 is retained for reporting compatibility. Each run
exports per-user vectors before aggregation. The final training seeds are 42,
2027, and 31415. A missing or failed required seed makes the confirmatory
aggregate incomplete; it is not imputed or replaced after the test is opened.

For a model $m$ and locked comparator $b$, the primary contrast is the mean of
per-user differences across the three seed occurrences:
\begin{equation}
  \Delta_{\mathrm{NDCG}} =
  \operatorname{mean}_{s,u}\left[
  \mathrm{NDCG@10}_{m,s,u} - \mathrm{NDCG@10}_{b,s,u}\right].
\end{equation}
The planned interval is a two-sided percentile 95\% hierarchical paired
bootstrap with 2,000 replicates. Each replicate resamples seed occurrences and,
within each occurrence, resamples user indices with replacement. The same user
indices are used for the model and comparator. The aggregate analysis uses a
NumPy PCG64 generator seeded with 42. No seed average is formed before the
bootstrap.

The comparator, configuration, cohort builder, and evaluator are frozen before
test access. Exploratory multi-model comparisons report the number of contrasts
and apply Holm correction. Efficiency measures, including latency, throughput,
and peak memory, are opened only after the accuracy artifact passes its required
integrity checks. These choices keep statistical uncertainty separate from
implementation reproducibility.

\section{Experimental Design and Evidence Scope}

\subsection{Evidence namespaces}

The paper uses four explicitly separated evidence surfaces. The first,
\texttt{PUBLIC\_DATA\_PROTOCOL\_VALIDATION}, contains the completed Pop/BPR MovieLens
lane and supports descriptive statements about procedural execution and
artifact traceability. The second, \texttt{OFFICIAL\_PROTOCOL\_REPRODUCTION}, is closed
as incomparable and contributes no result row. The third,
\texttt{HARMONIZED\_V5\_COMPARISON}, is sealed with major limitations and contributes no
numeric row to this manuscript. The fourth is a future follow-up namespace for
the v5 experiment and requires a new packet, registry, and audit.

This separation is substantive rather than administrative. A public run can
validate an execution path while using a different target event, dataset,
feature vocabulary, or source protocol. Conversely, a controlled benchmark can
answer a tightly defined internal question without supporting a public-data
generalization. The manuscript therefore does not pool metrics across these
surfaces or use one surface to fill a missing row in another.

Admission is applied at the row level rather than to an entire project label.
For a public row, admission requires a resolvable dataset binding, protocol
binding, model configuration, checkpoint rule, evaluator receipt, and
aggregate artifact. The four admitted rows satisfy that descriptive chain.
They are retained because their provenance is replayable, not because their
metric values are large or because they answer the primary Vietnamese retail
question. The official source-native lane is retained as a separate
incomparable outcome when its center, split, or evaluator cannot be aligned
without changing the source contract. The harmonized v5 lane is likewise
retained as a sealed limitation when mandatory comparator or statistical
conditions are absent. These states are informative: they distinguish an
execution path that was auditable from a scientific comparison that was not
yet admissible. A later experiment may extend the evidence, but it must do so
under a new namespace and a new admission decision rather than overwrite one
of these states.

\subsection{Run and artifact sequence}

The follow-up execution begins with a dataset snapshot and protocol manifest.
The runner then validates the environment, materializes the split, fits each
model using its permitted training inputs, exports score chunks in the frozen
catalog order, and invokes the independent evaluator. The resulting chain is
\textit{run manifest} to configuration and environment receipts, dataset
binding, checkpoint manifest, score artifact, per-user metrics, aggregate
metrics, and bootstrap report. Every link carries a SHA-256 binding or a
replayable derivation from the preceding artifact.

The completed public lane follows the same evidence shape at a smaller scope.
Its aggregate receipt is bound to SHA-256
\texttt{512fb8ad8862f5e677ec18f369d310ae58c9d6f3c4f7b2e79c17ede389930c17}.
The lane contains one Pop run with seed 42 and three BPR runs with seeds 42,
2027, and 31415. The public result table is generated from that receipt and the
four admitted row bindings; it is not reconstructed from console output.

\subsection{Follow-up comparison and ablations}

The controlled experiment first runs validation-only smoke checks and locks the
comparator registry. It then trains the independent deep, rule-only, and full
Hybrid conditions under the same data split, candidate catalog, masking, and
evaluator. The full Hybrid can be interpreted only if its score artifact is
finite, its catalog order matches the protocol, and its component artifacts are
present.

Three seeds are run after the configuration and checkpoint rule are frozen. The
test set is opened once, followed by per-user metrics and the hierarchical
bootstrap. The rule mechanism is evaluated on the train-defined rule cohort,
not on a cohort selected after seeing test results. The item-cold analysis
reports strict zero-edge and sparse denominators separately. If the full model
does not improve the locked contrast, the result is reported as null or
negative; changing the metric or dataset after the observation is outside the
protocol.

\subsection{External validity and adaptation}

An external dataset is admissible only when it preserves the essential task and
input contract: persistent user meaning, purchase or basket relevance, item
content required by the model, an auditable split, compatible candidate
semantics, and an evaluator mapping. A session or next-interaction dataset may
be useful for sensitivity analysis, but it is not automatically a purchase
replication. Similarly, replacing a source encoder or adaptation procedure with
an available text model changes the method contract. Transferable recommender
architectures retain explicit assumptions about target-domain adaptation
\cite{hou2022_unisrec} %<!--ref:hou2022_unisrec--><!--anchor:section:Figure 1; Sections 2.2-2.4, especially 2.4 Adaptation to Downstream Recommendation Tasks-->
\cite{hou2023_vqrec} %<!--ref:hou2023_vqrec--><!--anchor:section:Sections 2.1-2.3; Section 3.1 Experimental Setup-->
\cite{zheng2026_utgrec} %<!--ref:zheng2026_utgrec--><!--anchor:section:Section 2.4.2, Downstream Fine-tuning; Sections 3.1.1 and 3.1.3-3.1.4; Section 3.3, Ablation Study-->.
An external result is therefore labeled according to its actual alignment and
deviations, not according to a shared retail application label.

\section{Results}

\subsection{PUBLIC\_DATA\_PROTOCOL\_VALIDATION}

The admitted public lane completed four run rows: one deterministic Pop run and
three BPR runs. Table~\ref{tab:public-validation} reports the five metrics
retained in the aggregate receipt. The rows are descriptive protocol evidence;
the lane has no paired per-user vectors or inferential test, and the single Pop
run does not provide a seed-level variance estimate.

\begin{table}[htbp]
\centering
\caption{Descriptive results in the \texttt{PUBLIC\_DATA\_PROTOCOL\_VALIDATION} namespace.}
\label{tab:public-validation}
\small
\begin{tabular}{llrrrrr}
\toprule
Model & Seed & Recall@10 & MRR@10 & NDCG@10 & Hit@10 & Precision@10 \\
\midrule
Pop & 42 & 0.042678 & 0.053217 & 0.031164 & 0.186837 & 0.021444 \\
BPR & 42 & 0.101423 & 0.105435 & 0.070481 & 0.283439 & 0.038535 \\
BPR & 2027 & 0.106237 & 0.113953 & 0.075893 & 0.299363 & 0.041507 \\
BPR & 31415 & 0.109880 & 0.116951 & 0.077937 & 0.301486 & 0.042569 \\
\bottomrule
\end{tabular}
\end{table}

For the primary descriptive metric, the three BPR values have mean 0.074770
and sample standard deviation 0.003853. The Pop point estimate is 0.031164.
The BPR point estimates are consistent in direction across the three recorded
seeds under this exact public protocol, but the design does not support an
inferential conclusion because the Pop arm has one run and the lane lacks the
paired per-user artifact required for such a comparison. The other metrics show
the same row-level pattern, with BPR means of 0.294763 for Hit@10, 0.112113 for
MRR@10, 0.105847 for Recall@10, and 0.040870 for Precision@10.

\subsection{Evidence interpretation}

The public rows demonstrate that the locked execution path can preserve a
dataset binding, an environment receipt, a full-sort result, a checkpoint rule,
and an immutable aggregate receipt across a small reference lane. They do not
measure the proposed Hybrid, the v5 item-cold cohort, the rule-aligned cohort,
or external commerce validity. The official reproduction lane remains
incomparable, so no source-native target is imported into this table. The
controlled v5 rows are likewise withheld from the result namespace used here.

The table should therefore be read as an audit sample of a protocol, not as a
ranked leaderboard. Its four rows represent one deterministic control and
three stochastic reference runs; they do not provide balanced replication
across all methods or seeds. The BPR mean and dispersion summarize the
recorded BPR runs only, while the Pop value is a single point observation.
Neither summary changes the unit of analysis or supplies a paired comparison.
In particular, a row-level difference between the two labels would combine
different run counts and would not be a valid estimate of the primary v5
contrast. Keeping this limitation beside the table prevents a descriptive
sanity check from acquiring the rhetorical role of a scientific baseline.

\section{Discussion and Limitations}

\subsection{What the current evidence supports}

The main result of this stage is a protocol boundary. A recommender comparison
becomes more interpretable when the prediction event, catalog, masking,
relevance, evaluator, and training conditions are recorded together. The
public lane provides a concrete, audited example of that boundary: four rows
are traceable to a fixed environment, dataset binding, configuration, and
checkpoint chain. The value of these rows is procedural. They show what was run
and how the output can be traced, not what a different model or dataset would
produce.

The design also clarifies why the proposed architecture is an empirical
question. A rule score and a deep score can be combined in a single ranking
function, but a combined output does not identify the contribution of either
component. The no-Wide and deep-only conditions, train-only rule fitting, and
shared evaluator are needed to distinguish a mechanism effect from differences
in inputs, sampling, or optimization. Prior hybrid systems and basket studies
support this component-oriented design, while retaining source-specific task
assumptions \cite{peng2022_ham} %<!--ref:peng2022_ham--><!--anchor:section:Modeling sections for long-term preference and association components; Section 6.6 ablation study and Table 13-->
\cite{peng2023_m2} %<!--ref:peng2023_m2--><!--anchor:section:Section 4.1 model components; Section 6.5 ablation study-->
\cite{li2023_repetition_exploration} %<!--ref:li2023_repetition_exploration--><!--anchor:section:Introduction definitions and research questions; exploration-only analyses-->.

The namespace policy has a second benefit. It makes an incomplete or
incomparable source reproduction visible without treating a convenient public
run as a replacement. It also keeps the controlled benchmark honest about its
status as a constructed Vietnamese retail setting. This distinction is
especially important for cold-item claims: graph methods learn from observed
edges, and content signals can carry popularity structure, so neither sparse
graph results nor available text alone identifies performance for a strict
zero-edge item \cite{he2020_lightgcn} %<!--ref:he2020_lightgcn--><!--anchor:section:Sections 3.1.1 Light Graph Convolution, 3.1.2 Layer Combination, and 3.3 Model Training-->
\cite{meehan2025_cold_popbias} %<!--ref:meehan2025_cold_popbias--><!--anchor:section:Section 1, Introduction; Sections 3.1-3.3-->.

This boundary also gives a practical interpretation of the current dataset
decision. A catalog with 5,200 items and 5,000 users is not invalid merely
because its interaction matrix is sparse. Sparse data can be appropriate for a
controlled mechanism study when the target events, eligible-user count,
item-coverage profile, and rule-support distribution are measured before the
test is opened. The relevant question is whether enough users and items remain
in each predeclared cohort to make the estimand stable and interpretable. The
follow-up snapshot must consequently report the number of users with a novel
positive, the number of strict zero-edge and sparse items, the number of
transactions available for rule mining, and the exclusions caused by missing
content or malformed timestamps. If those diagnostics show that the rule or
cold-item estimand is underpopulated, the appropriate response is a versioned
dataset amendment or a reduced-method ablation. It is not to change the
generator, cohort definition, or metric after observing the test output.

The staged evidence policy is useful beyond this particular benchmark. It
allows infrastructure and provenance to be validated on a public dataset
without claiming that public-data behavior transfers to a constructed local
setting. It also prevents a difficult source reproduction from being hidden by
an easier surrogate. For a methods paper, this separation makes a negative
or incomplete empirical outcome publishable as an explicit boundary rather
than forcing every lane into a single success/failure label.

\subsection{Limitations}

First, the current public evidence is not an evaluation of the proposed model.
It contains a deterministic sanity control and one learned reference family,
with three BPR seeds but only one Pop run. Aggregate metrics are available, but
paired per-user vectors, confidence intervals, and bootstrap estimates are not
available in this lane. The descriptive direction should therefore not be
treated as a general algorithm ranking.

Second, the strict source-native reproduction lane is closed as incomparable.
This prevents a paper number from being transferred across a changed split,
candidate policy, or evaluator, but it also means that the current manuscript
cannot report a successful source-center reproduction. The follow-up must either
resolve the exact source contract or retain the incomparable status.

Third, the v5 benchmark is controlled and semi-synthetic. It is useful for
testing a declared mechanism under a reproducible schema, but its generated
behavior may not represent the distribution, incentives, or repeat/explore
balance of a live retail population. Next-basket work shows why repeat and novel
item behavior should be separated rather than summarized by one aggregate
quantity \cite{li2023_nbr_reality} %<!--ref:li2023_nbr_reality--><!--anchor:section:Figure 1 and Introduction contribution bullets; repeat-ratio and baseline analyses-->
\cite{li2023_mask_swap} %<!--ref:li2023_mask_swap--><!--anchor:section:Sections 4.1 and 4.2 mask/swap strategy; Section 5 BTBR model; Section 6 experiments-->.
The manuscript consequently treats the v5 counts as manifest declarations until
the new snapshot lineage is admitted and does not describe the records as
observed shopper logs.

Fourth, cold-item and content conclusions remain narrow. A strict zero-edge
definition is not interchangeable with a rare-item threshold, and the
denominators used in other cold-start studies do not automatically apply here.
Content representations can be aligned to collaborative teachers and inherit
their popularity structure; a text encoder can therefore be useful without
settling the ranking question \cite{huang2023_aldi} %<!--ref:huang2023_aldi--><!--anchor:section:Section 2.1, Problem Formulation; Section 2.2, Existing Solutions-->
\cite{meehan2026_semco} %<!--ref:meehan2026_semco--><!--anchor:section:Section 1, Introduction; Section 3.1, Problem Definition; Sections 3.2-3.3-->.

Fifth, external validity is not established. A compatible public commerce
dataset must preserve user, purchase, item-content, split, and candidate
semantics. Shared authorship, platform, or benchmark lineage also limits the
number of independent replications that can be claimed. These restrictions are
particularly relevant for transfer methods whose downstream adaptation is part
of the method rather than an optional implementation detail.

Finally, the current stage does not establish online behavior, revenue, latency
gates, resource efficiency, or deployment readiness. The efficiency analysis is
intentionally deferred until accuracy artifacts pass their integrity checks.
These limitations are design constraints for the next experiment, not missing
values to be filled by importing source-native metrics.

\subsection{Follow-up decision criteria}

The next empirical packet can support a model-effect conclusion only when the
v5 snapshot, protocol, comparator registry, component artifacts, three seeds,
per-user metrics, and bootstrap report are all sealed. The rule branch must be
evaluated on its predeclared cohort, the cold analysis must retain separate
denominators, and any external run must document deviations from the source
contract. If any of these conditions fails, the corresponding conclusion stays
unresolved. If the Hybrid contrast is null or negative, that outcome is
reported directly and the paper's interpretation is narrowed accordingly.

\section{Conclusion}

This paper defines a reproducible protocol for comparing retail recommenders
under a fixed temporal, novel-purchase, full-catalog ranking task. Its central
contribution is to bind data, split, candidate catalog, masking, evaluator,
model artifacts, and uncertainty analysis into an auditable evidence chain. The
completed public validation lane contains four descriptive Pop/BPR rows and
shows that the chain can be executed and traced in a separate public namespace.

The evidence does not yet answer whether the proposed decoupled Wide-and-Deep
Hybrid shows a positive contrast on the controlled Vietnamese benchmark. That question requires
the separately audited v5 follow-up with component ablations, cold-item
denominators, per-user statistics, and an explicit external-compatibility gate.
Until those artifacts are available, the appropriate conclusion is a bounded
methods contribution and a transparent empirical agenda.
The next gate is therefore operationally precise: it must either produce a
complete, null, or negative comparison under the frozen contract, or document
why the relevant estimand remains unresolved. Both outcomes are more
informative than importing a result from a different protocol.

\section*{Data Availability}

The public validation uses a named public dataset under its provider terms. The
controlled v5 data are a versioned research asset whose release and
redistribution status must be checked independently; this manuscript does not
assert unrestricted redistribution of row-level histories.

\section*{Ethics Statement}

The reported work uses secondary public data and controlled/generated behavior;
it does not recruit or interact with human participants. Any future use of
private logs, identifiable records, surveys, or online experiments requires a
new institutional ethics and data-protection determination.

\section*{Author Contributions}

Author metadata and contribution roles are intentionally left blank in this
venue-neutral version.

\section*{Funding and Conflicts of Interest}

Funding information and conflict-of-interest disclosures will be supplied with
the author and submission metadata.

\bibliographystyle{IEEEtran}
\bibliography{refs_stage2_5}

\end{document}
